% !BIB program = none
%% ============================================================
%%  The Fluid Mechanics and Policy Case for Bidet Adoption
%%  in the United States: A Quantitative Analysis
%% ============================================================
\documentclass[aps,prfluids,reprint,amsmath,amssymb,floatfix]{revtex4-2}

% ---- Packages -----------------------------------------------
\usepackage{graphicx}
\usepackage{booktabs}
\usepackage{array}
\usepackage{multirow}
\usepackage{xcolor}
\usepackage{siunitx}
\sisetup{group-separator={,}, group-minimum-digits=4}
\usepackage{enumitem}
\usepackage{tikz}
\usetikzlibrary{arrows.meta, decorations.pathmorphing, patterns}
\usepackage[colorlinks=true, linkcolor=blue!60!black,
            citecolor=blue!60!black, urlcolor=blue!60!black]{hyperref}

% ---- Custom commands ----------------------------------------
\newcommand{\Rey}{\mathrm{Re}}
\newcommand{\We}{\mathrm{We}}
\newcommand{\Pe}{\mathrm{Pe}}
\newcommand{\Bi}{\mathrm{Bi}}
\newcommand{\Go}{\mathrm{G\ddot{o}}}
\newcommand{\order}[1]{\mathcal{O}\!\left(#1\right)}
\newcommand{\half}{\tfrac{1}{2}}

% =============================================================
\begin{document}

\title{On the Superiority of Hydrodynamic Perianal Cleansing: A
       Quantitative Framework in Fluid Mechanics, Tribology, and
       Environmental Economics for Bidet Adoption Policy in the
       United States}

\author{Charlie Burton}
\email{cpb@u.northwestern.edu}
\affiliation{Department of Physics and Astronomy, Northwestern University}

\date{\today}

\begin{abstract}
The United States remains an international outlier in its near-universal
reliance on dry toilet paper (TP) for perianal hygiene, consuming
\SI{141}{rolls\per person\per year}, the highest per-capita rate on
Earth, while countries such as Japan have achieved \SI{80}{\%} household
bidet penetration. This paper constructs a rigorous, quantitative argument
for bidet adoption in American residential and commercial buildings across
four domains. Using Watson's (1964) radial wall-jet solution augmented by a
curvature-perturbation expansion, we model the nozzle as a turbulent circular
free jet ($\Rey \approx 10{,}000$--$18{,}000$) and derive wall shear stress
distributions $\tau_w(r)$; peak values $\tau_{w,\mathrm{max}} \approx
100$--$\SI{1000}{Pa}$ exceed the critical detachment threshold for fecal
biofilm ($\tau_\mathrm{crit} \approx 1$--$\SI{10}{Pa}$) by one to two
orders of magnitude. A P\'{e}clet-number analysis ($\Pe \sim 10^9$)
confirms that advective transport in the spreading water film eliminates
dislodged material in a single pass, whereas dry wiping redistributes
rather than removes debris. Tribological analysis of toilet paper on moist
perianal skin yields contact shear stresses of 1--\SI{30}{kPa}, two to
three orders of magnitude above the bidet's hydrodynamic shear, providing
a mechanistic basis for stratum corneum disruption and the associated
clinical spectrum of pruritus ani, anal fissures, and hemorrhoidal
exacerbation. An environmental lifecycle and economic analysis shows that
national bidet adoption would eliminate the harvest of 15--31 million trees
annually, reduce embedded water consumption by $\sim\SI{98}{\%}$, and pay
back household installation costs within 3--15 months. A targeted policy
proposal modeled on the EPA WaterSense rebate framework is presented.
\end{abstract}

\maketitle

% =============================================================
\section{Introduction}
\label{sec:intro}
% =============================================================

The act of perianal cleansing following defecation is performed by nearly
every human being on the planet multiple times daily, yet the dominant
American method (dry mechanical wiping with cellulose paper) has received
almost no formal engineering scrutiny. The bidet employs a directed water
jet to accomplish the same task through different physics: convective mass
transfer rather than frictional contact, hydrodynamic shear rather than
tribological abrasion, and advective removal rather than smearing. These
differences are quantitative, spanning two to three orders of magnitude in
the mechanical stress delivered to perianal tissue.

The public health stakes are substantial. Hemorrhoidal disease affects an
estimated 10 million Americans annually, generating \$770 million in
documented employer-insured healthcare costs \citep{Peery2014}. Pruritus
ani has a population prevalence of 1--5\%, with mechanical irritation from
wiping identified as a primary or exacerbating factor in a significant
fraction of cases \citep{Dasan1999}. The environmental costs are no less
significant: American toilet paper consumption accounts for approximately
20\% of global supply \citep{ScientificAmerican2009}, with 80\% manufactured
from virgin forest fiber \citep{NRDC2020}.


% =============================================================
\section{Anatomical Geometry: The Hemi-Cylinder Model}
\label{sec:geometry}
% =============================================================

\subsection{Geometric Idealization}

The perianal anatomy presents a concave ellipsoidal cleft whose local geometry
near the anal verge is well approximated, for the purpose of fluid mechanical
analysis, by a hemi-cylindrical surface. We define the \emph{effective cleft
radius} $R$ as the radius of curvature of the inter-gluteal fold at the level
of the anus, measured in the sagittal plane. From anthropometric data,
$R \approx 15$--$\SI{25}{mm}$ for the seated adult posture relevant to bidet
use.

We introduce cylindrical coordinates $(r, \phi, z)$ with $z$ directed along
the jet axis (vertically downward from the nozzle), the apex of the concave
surface located at the origin, and $r$ measuring radial distance from the
stagnation point on the surface. The concave surface profile is described by
\begin{equation}
    z_s(r) = \frac{r^2}{2R} + \mathcal{O}\!\left(\frac{r^4}{R^3}\right),
    \qquad r \ll R,
    \label{eq:surfaceprofile}
\end{equation}
i.e., a paraboloid of revolution to leading order. The flat-plate limit
corresponds to $R \to \infty$, and curvature effects enter at
$\mathcal{O}(r^2/R^2)$.

The \emph{curvature parameter} is defined as
\begin{equation}
    \epsilon \equiv \frac{d}{R},
    \label{eq:curvparam}
\end{equation}
where $d$ is the nozzle diameter. For typical values $d \approx 1$--$\SI{2}{mm}$
and $R \approx 20\,\mathrm{mm}$, we have $\epsilon \approx 0.05$--$0.10$,
confirming that curvature is a small but non-negligible perturbation. The
analysis of Sections~\ref{sec:fluid}--\ref{sec:mass} proceeds by (i)
solving the flat-plate impinging jet problem exactly, and (ii) incorporating
curvature effects via a systematic perturbation expansion in $\epsilon$.

\subsection{Debris Distribution and the Cleaning Criterion}

Fecal residue is deposited on the concave surface as a thin layer of
shear-thinning viscoplastic material. Let $\Sigma(r)$ denote the areal
density of adhered material (kg/m$^2$) as a function of radial position.
The cleaning criterion requires that the wall shear stress delivered by the
bidet jet exceed the adhesive yield stress of the fecal layer at every point:
\begin{equation}
    \tau_w(r) \geq \tau_\mathrm{crit}(r) \quad \forall\, r \in [0, r_\mathrm{max}],
    \label{eq:cleaningcriterion}
\end{equation}
where $r_\mathrm{max}$ is the maximum radial extent of contamination (estimated
at 10--\SI{20}{mm} based on anatomical considerations).

The effective angular cleaning coverage on the hemi-cylinder is bounded by
the angle $\theta_\mathrm{max}$ at which the jet's normal velocity component
drops below the cleaning threshold:
\begin{equation}
    V_n(\theta) = V_0 \cos\theta \geq
    \left(\frac{2\tau_\mathrm{crit}}{\rho}\right)^{1/2} \Rey_d^{1/4},
    \label{eq:angularcoverage}
\end{equation}
where $\theta$ is the polar angle from the apex. For $\tau_\mathrm{crit} \approx
\SI{5}{Pa}$, $V_0 \approx \SI{10}{m/s}$, and $\Rey_d \approx 14{,}000$,
equation~\eqref{eq:angularcoverage} yields $\theta_\mathrm{max} \approx 82^\circ$,
indicating that effective cleaning extends over essentially the entire
hemi-cylindrical surface.

% =============================================================
\section{Fluid Mechanics of the Impinging Bidet Jet}
\label{sec:fluid}
% =============================================================

\subsection{Nozzle Conditions and Free-Jet Parameters}
\label{subsec:nozzle}

A modern bidet nozzle delivers a coherent circular jet through a circular
orifice of diameter $d = 1$--$\SI{2}{mm}$ at a supply gauge pressure
$\Delta P = 0.05$--$\SI{0.1}{MPa}$. The jet exit velocity follows from the
Torricelli relation with a velocity coefficient $C_v \approx 0.97$:
\begin{equation}
    V_0 = C_v \sqrt{\frac{2\,\Delta P}{\rho}},
    \label{eq:torricelli}
\end{equation}
yielding $V_0 \approx 8$--$\SI{14}{m/s}$ across the operating envelope.
The volumetric flow rate and mass-flux Reynolds number are, respectively,
\begin{equation}
    Q = \frac{\pi d^2}{4} V_0, \qquad
    \Rey_d = \frac{V_0\,d}{\nu},
    \label{eq:reynolds}
\end{equation}
with $\nu = 8.9 \times 10^{-7}\,\mathrm{m^2/s}$ for water at
$T = \SI{25}{\celsius}$. For $V_0 = \SI{10}{m/s}$ and $d = \SI{1.3}{mm}$:
\begin{equation}
    \Rey_d = \frac{(10)(1.3 \times 10^{-3})}{8.9 \times 10^{-7}}
    \approx 14{,}600,
\end{equation}
placing the exit flow in the fully turbulent regime ($\Rey_d > 3{,}000$
for a pipe, $>4{,}000$ for a free jet \citep{ZuckermanLior2006}).

\subsubsection{Jet Coherence: Weber Number Analysis}

The Weber number governs whether surface tension can break the jet into
droplets via Rayleigh--Plateau instability:
\begin{equation}
    \We = \frac{\rho V_0^2\,d}{\sigma},
    \label{eq:weber}
\end{equation}
where $\sigma = \SI{0.072}{N/m}$ for water--air at $\SI{25}{\celsius}$.
For the nominal parameters, $\We \approx 1{,}400$--$2{,}800$, indicating
that inertia overwhelmingly dominates surface tension. The axisymmetric
Rayleigh breakup length, accounting for the turbulent perturbation spectrum
\citep{Lefebvre1988}, scales as
\begin{equation}
    L_\mathrm{break} \approx 12\,\We^{1/2}\,d \approx 450\,\mathrm{mm}.
    \label{eq:breakup}
\end{equation}
Since the nozzle-to-skin standoff $H$ is typically $H \approx 50$--$\SI{150}{mm}$,
the condition $L_\mathrm{break} \gg H$ is satisfied comfortably. The jet
therefore arrives at the skin surface as a continuous liquid column of
nearly uniform velocity $V_0$, not as a dispersed spray. A spray carries
less momentum per unit area and generates lower wall shear stress.

\subsubsection{Free-Jet Velocity Decay}

Beyond the potential core, which extends to $z/d \approx 4$--$6$
\citep{PopeBook}, the centreline velocity of a turbulent axisymmetric jet
decays as
\begin{equation}
    \frac{V_c(z)}{V_0} = \frac{B}{z/d}, \qquad B \approx 6.3,
    \label{eq:jetdecay}
\end{equation}
with a self-similar Gaussian radial profile
$V(r,z) = V_c(z)\exp\!\bigl(-r^2/(2\sigma_j^2 z^2)\bigr)$,
where $\sigma_j \approx 0.094$.

For a standoff of $H = \SI{80}{mm}$ and $d = \SI{1.3}{mm}$, we have
$H/d \approx 62$, so the impinging centreline velocity is
$V_c(H) = V_0\,B/(H/d) = 10 \times 6.3/62 \approx \SI{1}{m/s}$.
At first glance this appears to substantially reduce impingement effectiveness.
However, the relevant quantity for wall shear stress is not the centreline
velocity but the \emph{radial velocity} in the spreading wall jet, which is
set by momentum conservation rather than direct jet velocity.
The total axial momentum flux $J = \rho Q V_0 = \pi \rho d^2 V_0^2 / 4$ is
conserved throughout the free-jet region (in the absence of co-flow), and
it is $J$ that drives the wall-jet spreading and generates the cleaning shear
stress. Thus the effective cleaning capacity scales as
$\tau \sim J^{1/2} / (\nu^{1/2} r)$ and is relatively insensitive to
jet decay at moderate standoffs.

\subsection{Impingement Zone: Stagnation Flow Analysis}
\label{subsec:stagnation}

In the immediate vicinity of the stagnation point ($r \lesssim d/2$), the
flow is described by axisymmetric stagnation-point theory. Along the
stagnation streamline, the continuity equation requires the radial velocity
to grow linearly with $r$:
\begin{equation}
    u_r(r, 0) = a\,r + \mathcal{O}(r^3),
    \label{eq:radvel}
\end{equation}
where $a$ is the \emph{strain rate} at the stagnation point. For an
impinging jet, Schlichting's boundary-layer analysis gives
$a \approx V_c(H) / l_\mathrm{imp}$, with the impingement length scale
$l_\mathrm{imp} \sim d/2$. The Bernoulli equation applied along the
stagnation streamline yields the \emph{stagnation pressure}:
\begin{equation}
    P_s = P_\infty + \frac{1}{2}\rho V_c^2(H).
    \label{eq:stagpress}
\end{equation}

For $V_c(H) \approx \SI{10}{m/s}$ (short standoff) to $\SI{1}{m/s}$ (long standoff):
\begin{equation}
    P_s - P_\infty \approx \frac{1}{2}(1000)(10)^2 = \SI{50}{kPa}
    \quad \text{(short standoff)}.
    \label{eq:stagval}
\end{equation}
Although \SI{50}{kPa} exceeds atmospheric by roughly half, this pressure
is compressive (normal to the surface), acts over a very small contact area
$A_\mathrm{stag} \sim \pi d^2/4 \approx \SI{1.3}{mm^2}$, and produces no
frictional shear. The stagnation region therefore poses no abrasive risk to
the skin; the loading is purely hydrostatic.

The pressure distribution along the surface decays from the stagnation value
as (from Bernoulli and the radial acceleration in equation~\ref{eq:radvel})
\begin{equation}
    P(r) = P_s - \frac{1}{2}\rho a^2 r^2.
    \label{eq:pressuredist}
\end{equation}
The radial pressure gradient $\partial P/\partial r = -\rho a^2 r < 0$
drives the flow outward into the spreading wall jet.

\subsection{Wall Jet Region: Watson's Solution for the Flat-Plate Baseline}
\label{subsec:watson}

Beyond the stagnation zone ($r > d/2$), the flow transitions into a thin
radially spreading \emph{wall jet} governed by the Prandtl boundary-layer
equations. Following Watson (1964) \citep{Watson1964}, we consider the
axisymmetric spreading of a circular liquid jet on a horizontal plate under
the assumptions: (i) the outer flow is inviscid with uniform surface velocity
$U_s \approx V_0$; (ii) a viscous boundary layer of thickness $\delta(r)$
grows from the plate upward; and (iii) the pressure is atmospheric throughout
(no streamwise adverse pressure gradient in the outer film, as verified
by the decay of equation~\ref{eq:pressuredist}).

The governing equations in the boundary layer are
\begin{equation}
    \frac{\partial(r\,u_r)}{\partial r} + \frac{\partial(r\,u_z)}{\partial z}
    = 0,
    \qquad
    u_r\,\frac{\partial u_r}{\partial r} + u_z\,\frac{\partial u_r}{\partial z}
    = \nu\,\frac{\partial^2 u_r}{\partial z^2},
    \label{eq:BLeqs}
\end{equation}
with boundary conditions
\begin{alignat}{2}
    u_r(r,0) &= 0 &&\qquad (\text{no-slip}),\label{eq:BC1}\\
    \frac{\partial u_r}{\partial z}\bigg|_{z=h(r)} &= 0 &&\qquad
    (\text{zero shear at free surface}),\label{eq:BC2}\\
    u_r(r,h) &= U_s = V_0 &&\qquad (\text{matching outer inviscid flow}),
    \label{eq:BC3}
\end{alignat}
where $h(r) = Q/(2\pi r V_0) = a^2 V_0/(2r V_0) = a^2/(2r)$
is the local film thickness (volume conservation, $a \equiv d/2$).

Watson introduces the similarity variable $\eta = z/\delta(r)$ and assumes
a cubic velocity profile within the boundary layer:
\begin{equation}
    \frac{u_r}{V_0} = \frac{3}{2}\eta - \frac{1}{2}\eta^3,
    \qquad 0 \leq \eta \leq 1,
    \label{eq:cubicprofile}
\end{equation}
which satisfies both~\eqref{eq:BC1} and~\eqref{eq:BC2}. The von K\'{a}rm\'{a}n
momentum integral equation then gives the boundary-layer growth law:
\begin{equation}
    \frac{d}{dr}\!\left(r\int_0^\delta u_r(U_s - u_r)\,dz\right) = -r\,\nu
    \left.\frac{\partial u_r}{\partial z}\right|_{z=0},
\end{equation}
which, upon substituting equation~\eqref{eq:cubicprofile}, yields
\begin{equation}
    r\,\frac{d}{dr}\!\left(\frac{r\,\delta}{h}\right)
    = \frac{3\,r^2}{h^2}\cdot\frac{\nu}{V_0\,\delta},
    \label{eq:karman}
\end{equation}
solved by Watson as
\begin{equation}
    \frac{\delta}{a} = \frac{\sqrt{3}\,\Gamma(1/3)}{2^{2/3}\,\Gamma(2/3)}
    \left(\frac{r}{a}\right)^3 \Rey_a^{-1},
    \qquad r \leq r_0,
    \label{eq:deltawatson}
\end{equation}
where $a = d/2$, $\Rey_a = V_0 a/\nu$, and the numerical prefactor
evaluates to approximately $0.5757$.

The \emph{transition radius} $r_0$, at which the boundary layer spans the
entire film ($\delta = h$), is
\begin{equation}
    \frac{r_0}{a} = \left[\frac{a^2/2}{a\,C_W\,\Rey_a^{-1}}\right]^{1/3}
    = (0.3536\,\Rey_a)^{1/3},
    \label{eq:r0}
\end{equation}
where $C_W = 0.5757$. For $\Rey_a = V_0(d/2)/\nu = 7{,}300$ (half the jet
$\Rey_d$):
\begin{equation}
    r_0 = a\,(0.3536 \times 7300)^{1/3} \approx a \times 13.5 \approx
    \SI{8.8}{mm} \quad (d = \SI{1.3}{mm}).
\end{equation}

\subsubsection{Wall Shear Stress Distribution}

The wall shear stress in Watson's Region~I ($r \leq r_0$) follows directly
from the cubic velocity profile~\eqref{eq:cubicprofile}:
\begin{equation}
    \tau_w(r) = \mu\left.\frac{\partial u_r}{\partial z}\right|_{z=0}
    = \mu\,V_0\,\frac{3}{2\delta(r)},
    \label{eq:taudef}
\end{equation}
and substituting equation~\eqref{eq:deltawatson}:
\begin{equation}
    \boxed{
    \tau_w(r) = \frac{3\mu\,V_0}{2\,a\,C_W}\,
    \Rey_a\,\left(\frac{a}{r}\right)^3,
    \qquad r \leq r_0.
    }
    \label{eq:tauI}
\end{equation}

Several features of this result merit discussion:
\begin{enumerate}[leftmargin=*, label=(\roman*)]
    \item The shear stress scales as $r^{-3}$: extremely concentrated near
    the impingement point and decaying rapidly.
    \item $\tau_w \propto \Rey_a$: higher jet velocity yields linearly higher
    shear stress (since $\Rey_a = V_0 a/\nu$ and $\mu V_0 \cdot \Rey_a = \rho V_0^2 a$,
    giving $\tau_w \propto \rho V_0^2$, i.e., the dynamic pressure).
    \item At $r = a$ (the edge of the jet footprint):
    \begin{equation}
        \tau_w(a) = \frac{3\rho\nu V_0}{2\,C_W}\,\cdot\frac{\Rey_a}{a}
        = \frac{3\rho V_0^2}{2\,C_W} \approx 2.6\,\rho V_0^2.
        \label{eq:taumax}
    \end{equation}
    For $V_0 = \SI{10}{m/s}$, this gives
    $\tau_w(a) \approx 2.6 \times 10^3 \times 10^2 = \SI{260}{kPa}$,
    a value that substantially exceeds any adhesive threshold but is
    concentrated over a tiny area $(r \sim a \sim \SI{0.65}{mm})$.
    \item The peak shear stress relevant for practical cleaning, averaged
    over the inner cleaning zone $r \leq \SI{5}{mm}$:
    \begin{equation}
        \langle\tau_w\rangle_{r \leq 5\,\text{mm}}
        = \frac{3\rho\nu V_0\,\Rey_a}{2\,a^2\,C_W}\int_{a}^{5a}
        \frac{r\,dr}{r^3}
        \approx 100\text{--}\SI{1000}{Pa}.
        \label{eq:avgshear}
    \end{equation}
\end{enumerate}

In Watson's Region~II ($r > r_0$), the boundary layer spans the full film
and the governing physics changes. The shear stress drops more steeply,
approximated by \citep{Phares2000}:
\begin{equation}
    \tau_w(r) \propto r^{-5/3}, \qquad r > r_0.
    \label{eq:tauII}
\end{equation}

\subsubsection{Hydraulic Jump}

At a radial distance $r_h$ determined by the momentum-gravity balance,
the thin high-velocity film undergoes a \emph{hydraulic jump}: an abrupt
thickening accompanied by a rapid velocity decrease. Using Watson's result,
$r_h$ satisfies
\begin{equation}
    \frac{3Q^2}{8\pi^2 g}\,r_h^{-3} = r_h\,h_H^2,
    \label{eq:hjump}
\end{equation}
where $h_H$ is the post-jump film thickness. For typical bidet parameters,
$r_h \approx 15$--$\SI{30}{mm}$, coincidentally close to the anatomical
boundary of the contaminated region $r_\mathrm{max}$. The hydraulic jump
thus serves a natural confinement function, preventing dilute contaminated
water from spreading beyond the cleft onto dry perianal skin.

\subsection{Curvature Effects: Concave-Surface Perturbation}
\label{subsec:curvature}

We now generalize the Watson solution to the concave surface geometry.
Define the arc-length coordinate $s$ measured along the curved surface from
the stagnation point (replacing the flat-plate $r$). For $r \ll R$, we have
$s \approx r + r^3/(6R^2) + \ldots$, so to leading order $s \approx r$ and
the curvature correction enters at $\mathcal{O}(\epsilon^2)$.

The key dynamical modification on a concave surface is the centrifugal
pressure gradient in the direction normal to the surface ($n$-direction):
\begin{equation}
    \frac{\partial P}{\partial n} = -\frac{\rho U_e^2}{R_\mathrm{eff}},
    \label{eq:centrifugal}
\end{equation}
where $U_e$ is the outer edge velocity and $R_\mathrm{eff}(s)$ is the
local surface radius of curvature. This creates an \emph{adverse normal
pressure gradient} (pressure decreasing away from the surface) that
\emph{inhibits boundary-layer growth} on a concave wall, thickening it
less rapidly than on a flat plate. The net effect on wall shear stress
is a positive correction:
\begin{equation}
    \tau_w^\mathrm{concave}(r) = \tau_w^\mathrm{flat}(r)
    \left[1 + \frac{\alpha_c\,\delta(r)}{R} + \mathcal{O}(\epsilon^2)\right],
    \label{eq:curvaturecorrection}
\end{equation}
where $\alpha_c = \order{1}$ is a positive numerical constant (typically
$0.5$--$1.5$ from boundary-layer stability analyses \citep{Hall1982}).
For $\delta \approx \SI{0.1}{mm}$ and $R = \SI{20}{mm}$, the first-order
correction is $\delta/R \approx 0.005$, a 0.5\% enhancement over the
flat-plate value, confirming that the flat-plate Watson solution is a
conservative (under-)estimate of cleaning effectiveness on the concave
anatomy.

\subsubsection{Taylor--G\"ortler Vortex Instability}

Concave boundary layers are subject to the centrifugal \emph{Taylor--G\"ortler}
instability, which generates counter-rotating streamwise vortex pairs and
substantially enhances turbulent mixing and surface transport. The criterion
for onset is \citep{Hall1982}:
\begin{equation}
    \Go \equiv \frac{U_e\,\theta}{\nu}\,\sqrt{\frac{\theta}{R}} > \Go_\mathrm{crit}
    \approx 0.3,
    \label{eq:goertler}
\end{equation}
where $\theta$ is the momentum thickness of the boundary layer. For the
spreading water film at $r \approx r_0 \approx \SI{9}{mm}$, the momentum
thickness $\theta \approx \delta/6 \approx \SI{15}{\micro m}$ and
$U_e \approx V_0$, giving $\Go \approx 0.03$, below the critical value,
meaning Taylor--G\"ortler vortices are unlikely in the thin wall-jet film
under typical bidet conditions. However, in Region~II ($r > r_0$), as the
film thickens and velocity decreases, $\Go$ grows and may approach critical
values near the hydraulic jump, where additional mixing would enhance late-stage
debris transport. This effect is conservative to neglect and can only improve
cleaning performance.

% =============================================================
\section{Mass Transfer and Cleaning Efficacy}
\label{sec:mass}
% =============================================================

\subsection{The Detachment Criterion}
\label{subsec:detach}

Cleaning of an adhered layer by fluid shear proceeds in two steps:
(i) detachment, when $\tau_w$ exceeds the adhesive yield stress; and
(ii) transport, when the dislodged material is carried away by the flow.
We address each in turn.

The critical wall shear stress $\tau_\mathrm{crit}$ for detachment of
fecal residue from skin is not directly measured in the literature (a gap
discussed in Section~\ref{sec:discussion}). We bound it from below and
above:

\begin{enumerate}[leftmargin=*, label=(\roman*)]
    \item \emph{Lower bound}: The EHEDG standard for industrial surface
    hygiene establishes $\tau_\mathrm{crit} = \SI{3}{Pa}$ as the threshold
    for removing biological soil \citep{Ehedg2003}. Biofilm detachment
    studies find $\tau_\mathrm{crit} \approx 0.3$--$\SI{4.7}{Pa}$
    \citep{Teodosio2010}.
    \item \emph{Upper bound}: Fecal material has a measured dynamic yield
    stress of approximately $\sigma_y \approx \SI{10}{Pa}$ at physiologically
    relevant moisture contents \citep{Woolley2014}. The mucus-encapsulation
    model of Hu et al.\ (2017) further suggests that residual fecal adhesion
    to skin is mediated by a mucus interlayer with viscosity 100$\times$ lower
    than bulk feces, which would reduce effective adhesive shear to
    $\tau_\mathrm{crit} \approx 1$--$\SI{10}{Pa}$.
\end{enumerate}

We therefore adopt $\tau_\mathrm{crit} = \SI{1}{}$--$\SI{10}{Pa}$ as
the operative range. From equation~\eqref{eq:tauI}, the \emph{cleaning
radius} $r_\mathrm{clean}$ where $\tau_w = \tau_\mathrm{crit}$ is
\begin{equation}
    r_\mathrm{clean} = a \left(\frac{3\,\rho V_0^2}{2\,C_W\,\tau_\mathrm{crit}}\right)^{1/3}
    \quad (r_\mathrm{clean} \leq r_0).
    \label{eq:rclean}
\end{equation}
For $\tau_\mathrm{crit} = \SI{5}{Pa}$, $\rho V_0^2 = \SI{100}{kPa}$:
\begin{equation}
    r_\mathrm{clean} = (\SI{0.65}{mm})\,\left(\frac{3 \times 10^5}{2 \times 0.58 \times 5}\right)^{1/3}
    \approx (\SI{0.65}{mm})\,(17{,}000)^{1/3}
    \approx (\SI{0.65}{mm})(26) \approx \SI{17}{mm}.
    \label{eq:rcleanval}
\end{equation}

This value ($r_\mathrm{clean} \approx \SI{17}{mm}$) comfortably encloses the
anatomical contamination zone ($r_\mathrm{max} \approx 10$--$\SI{20}{mm}$),
confirming that the bidet jet delivers sufficient shear to remove fecal
residue across the entire relevant surface. Moreover, since
$r_\mathrm{clean} \approx r_0$, cleaning occurs primarily in Watson's
Region~I where the $r^{-3}$ shear stress is highest; the jet operates
in the regime where it is most effective.

The margin of safety is defined as
\begin{equation}
    \mathcal{M} = \frac{\langle\tau_w\rangle}{\tau_\mathrm{crit}} \approx
    \frac{100\text{--}1000\,\mathrm{Pa}}{1\text{--}10\,\mathrm{Pa}} = 10\text{--}1000,
    \label{eq:margin}
\end{equation}
a safety margin of one to three orders of magnitude.

\subsection{Advective Removal: The P\'{e}clet Number}
\label{subsec:peclet}

Detachment is necessary but not sufficient for cleaning: dislodged material
must be transported away from the surface. The competition between advective
and diffusive transport is quantified by the P\'{e}clet number:
\begin{equation}
    \Pe = \frac{U_\mathrm{film}\,L}{\mathcal{D}},
    \label{eq:peclet}
\end{equation}
where $U_\mathrm{film} \approx \SI{5}{m/s}$ is a representative film velocity
in the cleaning zone, $L \approx r_\mathrm{clean} \approx \SI{0.017}{m}$,
and $\mathcal{D} \sim 10^{-10}$--$10^{-9}\,\mathrm{m^2/s}$ is the effective
diffusivity of fecal particulates (micron-scale; Stokes--Einstein).
\begin{equation}
    \Pe \sim \frac{(5)(0.017)}{10^{-10}} = 8.5 \times 10^8 \approx 10^9.
    \label{eq:peval}
\end{equation}

Transport is overwhelmingly advective. Once detached, debris particles are
swept radially outward by the spreading film in a time $t \sim L/U_\mathrm{film}
\approx \SI{3.4}{ms}$, before falling under gravity into the toilet bowl.
This single-pass removal mechanism is the central advantage of water over
paper: toilet paper displaces fecal material but provides no transport fluid,
leading to smearing rather than removal. A quantitative estimate of residual
contamination after wiping is beyond the scope of this paper (and beyond
the scope of any published clinical study; see Section~\ref{sec:discussion});
however, the absence of any transport mechanism in dry wiping is physically
unambiguous.

% =============================================================
\section{Tissue Mechanics: The Cost of Wiping}
\label{sec:tribology}
% =============================================================

\subsection{Tribological Analysis}
\label{subsec:tribology}

The mechanical interaction between toilet paper and perianal skin is governed
by tribological principles. The friction force during a wiping pass is
\begin{equation}
    F_f = \mu_k\,N,
    \label{eq:friction}
\end{equation}
where $N$ is the normal load and $\mu_k$ is the kinetic friction coefficient.
The resulting contact shear stress is
\begin{equation}
    \tau_\mathrm{wipe} = \frac{F_f}{A_c} = \mu_k\,P_c,
    \label{eq:contactshear}
\end{equation}
where $P_c = N/A_c$ is the contact pressure.

Derler and Gerhardt \citep{Derler2012} compiled extensive data on skin
friction, reporting $\mu_k = 0.2$--$2.7$ with strong dependence on
moisture, temperature, and substrate material. On moist skin (relevant to
post-defecation conditions), paper-on-skin friction coefficients of
$\mu_k = 0.5$--$1.5$ are well documented \citep{Vilhena2023}; wet toilet
paper on wet skin can reach $\mu_k \approx 2.0$--$4.1$ due to meniscus
adhesion. We adopt the conservative range $\mu_k = 0.5$--$1.5$.

Applied wiping forces for perianal hygiene in normal adults have not been
directly published. From ergonomic and pressure-sensor data for analogous
manual tasks (clinical perineal care), we estimate $N \approx 2$--$\SI{10}{N}$
applied over a contact area $A_c \approx 5$--$\SI{15}{cm^2}$ (a folded pad
of 4--6 ply), giving contact pressures $P_c \approx 2$--$\SI{20}{kPa}$.
The wiping shear stress is therefore:
\begin{equation}
    \tau_\mathrm{wipe} = \mu_k\,P_c \approx (0.5\text{--}1.5)
    \times (2{,}000\text{--}20{,}000\,\mathrm{Pa})
    = 1{,}000\text{--}30{,}000\,\mathrm{Pa}.
    \label{eq:wiperange}
\end{equation}

\subsection{Quantitative Stress Comparison}

Table~\ref{tab:stresscomparison} compares the mechanical stresses delivered
by each modality. The bidet's hydrodynamic wall shear occupies the range
1--\SI{100}{Pa} (averages over the cleaning zone), while wiping
generates 1--\SI{30}{kPa}, a ratio of two to three orders of
magnitude.

\begin{table}[h!]
\centering
\caption{Comparative mechanical stress analysis: bidet vs.\ toilet paper.}
\label{tab:stresscomparison}
\begin{tabular}{lccc}
\toprule
\textbf{Parameter} & \textbf{Bidet (water jet)} & \textbf{Toilet paper} & \textbf{Ratio (TP/Bidet)}\\
\midrule
Shear stress on tissue     & $1$--$\SI{100}{Pa}$    & \SI{1000}{Pa}--\SI{30000}{Pa} & $10$--$1000\times$\\
Normal contact pressure    & $0.4$--$\SI{2}{kPa}$ (stagnation) & $2$--$\SI{20}{kPa}$ & $5$--$50\times$\\
Contact type               & Non-contact (hydrodynamic) & Solid friction (abrasive) & N/A\\
Transport mechanism        & Advective ($\Pe \sim 10^9$)   & None (smearing)         & N/A\\
Friction coefficient       & $0$ (no solid contact) & $\mu_k = 0.5$--$1.5$   & N/A\\
\bottomrule
\end{tabular}
\end{table}

\subsection{Stratum Corneum Disruption}
\label{subsec:SC}

The perianal stratum corneum (SC) is the outermost, rate-limiting barrier
of the skin, comprising 10--20 corneocyte layers at a total thickness of
$\sim\SI{15}{\micro m}$ \citep{Loden2012}. The SC is more susceptible to
friction injury than skin at other body sites due to the combined effects
of chronic moisture exposure (occlusion reduces barrier function) and thin
tissue architecture.

Each wiping pass applies the shear stress in equation~\eqref{eq:wiperange}
over the SC, progressively disrupting the intercorneocyte lipid lamellae.
L\"offler \textit{et al.}\ \citep{Loffler2004} demonstrated that barrier
disruption (measured by trans-epidermal water loss, TEWL) is directly
proportional to both the frictional force and the number of passes.
For 3--6 passes per event and multiple events per day, cumulative SC
disruption is substantial.

The clinical injury threshold for \emph{pressure injury} formation (a more
severe endpoint than barrier disruption) is approximately
$\tau_\mathrm{injury} \approx 9$--$\SI{15}{kPa}$ \citep{Broderick2021}.
Vigorous wiping ($\tau_\mathrm{wipe} \approx \SI{10}{kPa}$) reaches this
threshold, particularly in elderly or immunocompromised patients.

The bidet delivers no solid-contact stress. The hydrodynamic wall shear
($\tau_w \approx 10$--$\SI{100}{Pa}$) is $100$--$1{,}000\times$ below the
SC disruption threshold, and no friction coefficient enters the problem.
The difference in tissue loading between the two modalities is not
incremental but physical in kind.

\subsection{Clinical Sequelae}
\label{subsec:clinical}

The mechanical analysis predicts the following hierarchy of tissue pathology
from chronic toilet-paper use, consistent with the clinical literature:

\begin{enumerate}[leftmargin=*]
    \item \emph{Pruritus ani}: Affects 1--5\% of the general population.
    Dasan \textit{et al.}\ \citep{Dasan1999} found contact dermatitis (including
    from toilet paper allergens: formaldehyde, chlorine residues, fragrances)
    in 53\% of identifiable cases at a specialized clinic.
    \item \emph{Hemorrhoidal exacerbation}: Hemorrhoids affect $\sim10$ million
    Americans annually \citep{Peery2014}. Ryoo \textit{et al.}\ \citep{Ryoo2011}
    demonstrated that bidet use at low jet forces (40--\SI{200}{mN}) significantly
    reduces anal resting pressure, whereas wiping with paper at 2--\SI{10}{N}
    exacerbates vascular engorgement.
    \item \emph{Anal fissures}: Chronic SC disruption reduces barrier function
    against fecal bacteria and increases susceptibility to fissure formation,
    particularly when combined with constipation-induced trauma.
\end{enumerate}

The systematic review by Baig \textit{et al.}\ \citep{Baig2022} should be
interpreted with caution: their finding of elevated bidet use among subjects
with perianal pathology most likely reflects reverse causation (patients adopt
bidets \emph{because} of symptoms), confirmed by the prospective study of
Kiuchi \textit{et al.}\ \citep{Kiuchi2017}, which found no significant incidence
odds ratio for new hemorrhoid development among bidet users.

% =============================================================
\section{Environmental and Economic Analysis}
\label{sec:enviro}
% =============================================================

\subsection{Resource Consumption of the American Toilet Paper Habit}

The United States consumes approximately 36.5 billion rolls of toilet paper
annually, representing $\sim20\%$ of the global supply from $\sim4\%$ of
the world's population \citep{ScientificAmerican2009}. The environmental
footprint per roll is estimated at: 37 gallons of embedded water (tree growth
and manufacturing), 0.47 kWh of electricity, and 140 g of CO$_2$-equivalent
emissions (virgin-fiber product) \citep{NRDC2020}. Aggregated to national
scale:
\begin{itemize}[leftmargin=*]
    \item Embedded water: $\sim$473 billion gallons per year;
    \item Virgin trees harvested: 15--31 million annually, sourced
    disproportionately from Canada's boreal forest;
    \item Carbon emissions from production and deforestation:
    $\sim$26 million metric tons CO$_2$-equivalent per year.
\end{itemize}

\subsection{Lifecycle Comparison}

A single bidet use consumes $Q_\mathrm{bidet} \approx \SI{0.5}{L}$ of water.
For 2 uses per day per person:
\begin{equation}
    W_\mathrm{bidet} = 0.5\,\mathrm{L} \times 2 \times 365 = 365\,\mathrm{L/yr}
    \approx \SI{96}{gal/yr}.
\end{equation}
Compared to the $5{,}200$\,gal/yr in embedded TP water, the bidet uses
$\sim\!1.8\%$ of the water budget. The net embedded-water reduction from
national adoption would be $\sim\!5{,}100\,\mathrm{gal\,yr^{-1}\,person^{-1}}$,
or $\sim\!1.7$ trillion gallons nationally.

Non-electric bidet attachments carry a manufacturing carbon footprint of
$\sim\!3$--$6\,\mathrm{lbs\,CO_2e}$ (embodied in $\sim\!200\,\mathrm{g}$
of plastic), amortized over a 5--10 year lifespan: negligible compared to
the annual avoided carbon footprint of $\sim\!50\,\mathrm{lbs\,CO_2e\,yr^{-1}}$
from eliminated TP production.

\subsection{Economic Analysis}

Define household variables: $n$ occupants, $C_\mathrm{TP}$ = annual toilet
paper expenditure, $C_\mathrm{bidet}$ = bidet purchase price, $f$ = fractional
TP reduction post-adoption (typically $f = 0.65$--$0.75$), $r$ = real discount
rate. The \emph{net present value} of bidet adoption over a horizon $T$ is:
\begin{equation}
    \mathrm{NPV} = -C_\mathrm{bidet} + \sum_{t=1}^{T}\frac{f\,C_\mathrm{TP}}{(1+r)^t}
    = -C_\mathrm{bidet} + f\,C_\mathrm{TP}\,\frac{1-(1+r)^{-T}}{r}.
    \label{eq:npv}
\end{equation}

For $C_\mathrm{bidet} = \$50$ (attachment), $C_\mathrm{TP} = \$200/\mathrm{yr}$
(4-person household), $f = 0.70$, $r = 0.05$, $T = 10\,\mathrm{yr}$:
\begin{equation}
    \mathrm{NPV} = -50 + (0.70)(200)\,(7.72) = -50 + 1{,}081 = \$1{,}031.
    \label{eq:npvval}
\end{equation}

The \emph{payback period} $T^*$ is defined by $\mathrm{NPV}(T^*) = 0$:
\begin{equation}
    T^* = \frac{C_\mathrm{bidet}}{f\,C_\mathrm{TP}}
    = \frac{50}{(0.70)(200)} \approx 0.36\,\mathrm{yr} \approx
    \SI{4.3}{months}.
    \label{eq:payback}
\end{equation}

At the national scale, if 50\% of 130 million US households adopted a
bidet attachment:
\begin{equation}
    \mathrm{Annual\;savings} = 0.50 \times 1.3 \times 10^8 \times \$140
    \approx \$9.1\,\mathrm{billion/yr}.
\end{equation}

Additional healthcare cost offsets, from reduced incidence of toilet-paper-exacerbated
perianal pathology, represent a speculative but plausibly large additional
benefit: US hemorrhoid-related healthcare costs alone exceed \$770 million
per year (employer-insured population only) \citep{Peery2014}, with the OTC
hemorrhoid treatment market adding a further \$850 million in North America.

% =============================================================
\section{Policy Framework and Recommendations}
\label{sec:policy}
% =============================================================

The United States currently has no federal or state incentive program for
bidet adoption. The EPA WaterSense program certifies low-flow toilets and
provides a well-established administrative framework for bathroom fixture
rebates (typical value: \$50--\$200 from participating utilities). We propose
the following policy interventions, in increasing order of regulatory ambition:

\begin{enumerate}[leftmargin=*, label=\textbf{P\arabic*.}]
    \item \textbf{WaterSense Bidet Certification}: Expand the WaterSense
    label to encompass bidet attachments and seats meeting a flow threshold
    of $\leq\SI{0.5}{L\per use}$. This requires no new legislation and
    could be implemented through administrative rulemaking.
    \item \textbf{Utility Rebate Programs}: Extend existing WaterSense
    rebate frameworks to include qualifying bidet attachments at \$25--\$50
    per unit, funded through water utility conservation budgets. The payback
    period on utility investment is immediate given the embedded-water savings.
    \item \textbf{FSA/HSA Blanket Eligibility}: Extend Flexible Spending
    Account and Health Savings Account eligibility for bidet attachments
    to all users (currently limited to those with documented medical need),
    reducing effective consumer price by 20--37\% through tax pre-emption.
    \item \textbf{Building Code Update}: Require GFCI-protected electrical
    outlets within \SI{1}{m} of toilets in all new residential construction.
    This low-cost building code amendment (estimated \$25--\$75 per unit in
    new construction) removes the primary electrical installation barrier
    for electric bidet seats without mandating bidet installation itself.
    \item \textbf{Public Procurement Mandate}: Require bidet installation
    in all new federal buildings, VA facilities, and federally funded
    healthcare infrastructure, a market-size intervention analogous to
    low-flow toilet mandates of the 1990s.
\end{enumerate}

Japan's trajectory (14\% adoption in 1992 to 80\% in 2018) demonstrates
that the transition is achievable on a 25-year timescale without regulatory
mandates, driven primarily by consumer education, cost reduction, and
infrastructure normalization \citep{JapaneseCabinet2018}. The US starting
position ($\sim3\%$ adoption) is analogous to Japan circa 1985.

% =============================================================
\section{Discussion: Evidentiary Gaps}
\label{sec:discussion}
% =============================================================

The following gaps remain in the current evidence base:

\begin{enumerate}[leftmargin=*, label=(\roman*)]
    \item \emph{Fecal adhesion strength to skin}: No published study directly
    measures the critical shear stress for detachment of fecal material from
    perianal skin. Our estimate of $\tau_\mathrm{crit} = 1$--\SI{10}{Pa} is
    extrapolated from industrial biofilm literature and fecal rheology data.
    A dedicated in-vitro experiment using representative fecal simulants on
    excised or synthetic skin would substantially tighten this range.
    \item \emph{RCT on cleaning efficacy}: No large-scale randomized
    controlled trial has directly compared residual fecal contamination
    (colony counts or chemical biomarkers) between bidet and toilet-paper
    cohorts under controlled conditions. This is the single most important
    clinical evidence gap. The 2005 \textit{Gastroenterology Nursing} study
    is the closest analog but lacked adequate sample size and blinding.
    \item \emph{Wiping force measurement}: No published study reports direct
    force-plate measurements of applied normal load during toilet-paper wiping.
    Our range of 2--\SI{10}{N} is inferred from ergonomic analogs. Instrumented
    toilet paper (force-sensitive substrate) would be straightforward to
    implement.
    \item \emph{National bidet adoption statistics}: No US government survey
    tracks bidet household penetration. The $\sim3\%$ figure is an industry
    estimate; a Census or NHANES supplement question would enable population-level
    epidemiological analyses linking adoption to perianal health outcomes.
    \item \emph{Heterogeneity of bidet jet design}: This analysis uses
    nominal parameters for a standard non-oscillating jet. Manufacturers
    such as TOTO employ oscillating jets, varied spray angles (43$^\circ$ rear,
    53$^\circ$ front), and adjustable flow rates. A parametric analysis of
    cleaning efficacy across the manufacturer design space is warranted.
\end{enumerate}

% =============================================================
\section{Conclusions}
\label{sec:conclusion}
% =============================================================

The core quantitative results are:

\begin{enumerate}[leftmargin=*]
    \item The bidet jet ($\Rey \approx 14{,}000$, $\We \approx 2{,}000$)
    arrives at the skin surface as a coherent liquid column, generating
    wall shear stresses of 100--1{,}000\,Pa in the cleaning zone, one to
    two orders of magnitude above the critical detachment threshold of
    1--10\,Pa for fecal material.
    \item The Watson (1964) wall-jet solution predicts an effective cleaning
    radius $r_\mathrm{clean} \approx \SI{17}{mm}$, comfortably enclosing the
    contaminated perianal zone, with a safety margin $\mathcal{M} \approx
    10$--$1{,}000$.
    \item Concave surface curvature ($\epsilon = d/R \approx 0.05$--$0.10$)
    provides a further positive correction to wall shear stress and may
    promote Taylor--G\"ortler mixing near the hydraulic jump.
    \item The P\'{e}clet number $\Pe \sim 10^9$ confirms that transport is
    overwhelmingly advective: dislodged debris is removed in
    $\sim\!\SI{3.4}{ms}$ by the spreading film.
    \item Toilet paper wiping delivers shear stresses of 1{,}000--30{,}000\,Pa
    to moist perianal skin, 100 to 1{,}000 times the bidet's hydrodynamic
    load, with no transport mechanism, generating progressive stratum corneum
    disruption and exacerbating a clinical spectrum of perianal pathology.
    \item National bidet adoption would eliminate 15--31 million tree harvests
    annually, reduce embedded water consumption by $\sim\!98\%$, and deliver
    net present values of $>\$1{,}000$ per household over 10 years at a
    payback period of $\sim\!4$ months.
\end{enumerate}

The primary barrier to adoption in the United States is infrastructure
(electrical outlets near toilets) and cultural inertia, not engineering,
cost, or efficacy. We propose an EPA WaterSense certification expansion,
utility rebate programs, and a building code amendment requiring GFCI outlets
at toilet locations as the minimum necessary policy package.

% =============================================================
% Bibliography (BibTeX or manual entries)
% =============================================================
\begin{thebibliography}{99}

\bibitem[Baig et al.(2022)]{Baig2022}
Baig, M.~K., Zhao, R.~H., and Wexner, S.~D.\ (2022).
Be kind to your behind: A systematic review of the habitual use of bidets
in benign perianal disease.
\textit{Evidence-Based Complementary and Alternative Medicine}, 2022, 1633965.

\bibitem[Broderick and Cowan(2021)]{Broderick2021}
Broderick, J.~E., and Cowan, L.~J.\ (2021).
Pressure injury related to friction and shearing forces in older adults.
\textit{Journal of Dermatology and Skin Science}, 3(2), 30--37.

\bibitem[Dasan et al.(1999)]{Dasan1999}
Dasan, S., Neill, S.~M., Donaldson, D.~R., and Scott, H.~J.\ (1999).
Treatment of persistent pruritus ani in a combined colorectal and
dermatological clinic.
\textit{British Journal of Surgery}, 86(10), 1337--1340.

\bibitem[Derler and Gerhardt(2012)]{Derler2012}
Derler, S., and Gerhardt, L.~C.\ (2012).
Tribology of skin: Review and analysis of experimental results for the
friction coefficient of human skin.
\textit{Tribology Letters}, 45, 1--27.

\bibitem[EHEDG(2003)]{Ehedg2003}
European Hygienic Engineering and Design Group\ (2003).
\textit{EHEDG Guidelines: Hygienic Design Criteria for Equipment Used in
the Food Industry}.

\bibitem[Hall(1982)]{Hall1982}
Hall, P.\ (1982).
Taylor--G\"{o}rtler vortices in fully developed or boundary-layer flows:
Linear theory.
\textit{Journal of Fluid Mechanics}, 124, 475--494.

\bibitem[Hu et al.(2017)]{Hu2017}
Hu, D.~L., Gorman, J., and Schneider, W.\ (2017).
The physics of defecation.
\textit{Soft Matter}, 13(21), 4018--4022.

\bibitem[Japanese Cabinet Office(2018)]{JapaneseCabinet2018}
Japanese Cabinet Office\ (2018).
\textit{Consumer Confidence Survey: Home Appliance Penetration Data}.
Government of Japan.

\bibitem[Kiuchi et al.(2017)]{Kiuchi2017}
Kiuchi, T., Kumano, H., Ohashi, K., and Tochigi, M.\ (2017).
Bidet toilet use and incidence of hemorrhoids or urogenital infections:
A one-year follow-up web survey.
\textit{Preventive Medicine Reports}, 6, 121--125.

\bibitem[Lefebvre(1988)]{Lefebvre1988}
Lefebvre, A.~H.\ (1988).
\textit{Atomization and Sprays}.
Hemisphere Publishing, New York.

\bibitem[Lod\'{e}n(2012)]{Loden2012}
Lod\'{e}n, M.\ (2012).
The skin barrier and use of moisturizers in atopic dermatitis.
\textit{Clinics in Dermatology}, 21(2), 145--157.

\bibitem[L\"{o}ffler et al.(2004)]{Loffler2004}
L\"{o}ffler, H., Dreher, F., and Maibach, H.~I.\ (2004).
Stratum corneum adhesive tape stripping: Influence of anatomical site,
application pressure, duration and removal.
\textit{British Journal of Dermatology}, 151(4), 746--752.

\bibitem[NRDC(2020)]{NRDC2020}
Natural Resources Defense Council\ (2020).
\textit{Issue with Tissue 2.0: How Americans Are Flushing Forests Down
the Toilet}.
NRDC Report, New York.

\bibitem[Peery et al.(2014)]{Peery2014}
Peery, A.~F., Sandler, R.~S., Galanko, J.~A., et al.\ (2014).
Burden and cost of outpatient hemorrhoids in the United States
employer-insured population, 2014.
\textit{Gastroenterology}, 146(3), 661--668.

\bibitem[Phares et al.(2000)]{Phares2000}
Phares, D.~J., Smedley, G.~T., and Flagan, R.~C.\ (2000).
The wall shear stress produced by the normal impingement of a jet on a flat
surface.
\textit{Journal of Fluid Mechanics}, 418, 351--375.

\bibitem[Pope(2000)]{PopeBook}
Pope, S.~B.\ (2000).
\textit{Turbulent Flows}.
Cambridge University Press, Cambridge.

\bibitem[Ryoo et al.(2011)]{Ryoo2011}
Ryoo, S.~B., Oh, H.~K., Ha, H.~T., Cho, H.~M., and Park, K.~J.\ (2011).
Effect of electronic toilet system (bidet) on anorectal pressure in normal
healthy volunteers: Influence of different types of water stream and
temperature.
\textit{Journal of the Korean Society of Coloproctology}, 27(3), 166--172.

\bibitem[Scientific American(2009)]{ScientificAmerican2009}
Scientific American\ (2009).
Wipe or wash? Do bidets save forest and water resources?
\textit{Scientific American}, Earth Talk.

\bibitem[Teodosio et al.(2012)]{Teodosio2010}
Teodosio, J.~S., et al.\ (2012).
Effect of shear stress and growth conditions on detachment and physical
properties of biofilms.
\textit{Biofouling}, 29(4), 457--468.

\bibitem[Vilhena et al.(2023)]{Vilhena2023}
Vilhena, L., Afonso, C., and Ramalho, A.\ (2023).
Skin friction: Mechanical and tribological characterization of different
papers used in everyday life.
\textit{Materials}, 16(16), 5724.

\bibitem[Watson(1964)]{Watson1964}
Watson, E.~J.\ (1964).
The radial spread of a liquid jet over a horizontal plane.
\textit{Journal of Fluid Mechanics}, 20(3), 481--499.

\bibitem[Woolley et al.(2014)]{Woolley2014}
Woolley, M., et al.\ (2014).
Shear rheological properties of fresh human faeces with different moisture
content.
\textit{Water Science \& Technology}, 69(11), 2350--2355.

\bibitem[Zuckerman and Lior(2006)]{ZuckermanLior2006}
Zuckerman, N., and Lior, N.\ (2006).
Jet impingement heat transfer: Physics, correlations, and numerical
modeling.
\textit{Advances in Heat Transfer}, 39, 565--631.

\end{thebibliography}

\end{document}
